\BeleskeID{06-s014}
% element: 06-s014:d01
E predstavlja celobrojnu vrednost polja eksponenta, dok F u uslovima označava da li su svi biti frakcije nula. Oznaka NaN znači da zapis ne predstavlja brojnu vrednost, na primer pri matematički neodređenom rezultatu.

Kod subnormalnih brojeva vodeći bit je nula i faktor ostaje $2^{-126}$. Najmanji nenulti pozitivni zapis sa jedinicom u poslednjem bitu frakcije ima vrednost $2^{-126}2^{-23}=2^{-149}$. Tako se konačni brojevi približavaju nuli postepeno, bez skoka sa najmanjeg normalnog broja pravo na nulu.

Znak ostaje zapisan i za nulu i za beskonačnost. Zato postoje dva zapisa nule, sa S=0 i S=1, kao i zapisi pozitivne i negativne beskonačnosti.
